\documentclass[10pt,a4paper]{amsart}
\usepackage{amsmath,amssymb,amsthm,graphicx}
\usepackage{verbatim}
\newenvironment{reference}{\comment}{\endcomment}
%\newenvironment{reference}{}{}
\newenvironment{slogan}{\comment}{\endcomment}
\newenvironment{history}{\comment}{\endcomment}

% For commutative diagrams we use Xy-pic
\usepackage[all]{xy}

% We use 2cell for 2-commutative diagrams.
\xyoption{2cell}
% Theorem environments.
%
\theoremstyle{plain}
\newtheorem{theorem}[subsection]{Theorem}
\newtheorem{proposition}[subsection]{Proposition}
\newtheorem{lemma}[subsection]{Lemma}

\theoremstyle{definition}
\newtheorem{definition}[subsection]{Definition}
\newtheorem{example}[subsection]{Example}
\newtheorem{exercise}[subsection]{Exercise}
\newtheorem{situation}[subsection]{Situation}

\theoremstyle{remark}
\newtheorem{remark}[subsection]{Remark}
\newtheorem{remarks}[subsection]{Remarks}

\numberwithin{equation}{subsection}

% Macros
%
\def\lim{\mathop{\mathrm{lim}}\nolimits}
\def\colim{\mathop{\mathrm{colim}}\nolimits}
\def\Spec{\mathop{\mathrm{Spec}}}
\def\Hom{\mathop{\mathrm{Hom}}\nolimits}
\def\Ext{\mathop{\mathrm{Ext}}\nolimits}
\def\SheafHom{\mathop{\mathcal{H}\!\mathit{om}}\nolimits}
\def\SheafExt{\mathop{\mathcal{E}\!\mathit{xt}}\nolimits}
\def\Sch{\mathit{Sch}}
\def\Mor{\mathop{\mathrm{Mor}}\nolimits}
\def\Ob{\mathop{\mathrm{Ob}}\nolimits}
\def\Sh{\mathop{\mathit{Sh}}\nolimits}
\def\NL{\mathop{N\!L}\nolimits}
\def\CH{\mathop{\mathrm{CH}}\nolimits}
\def\proetale{{pro\text{-}\acute{e}tale}}
\def\etale{{\acute{e}tale}}
\def\QCoh{\mathit{QCoh}}
\def\Ker{\mathop{\mathrm{Ker}}}
\def\Im{\mathop{\mathrm{Im}}}
\def\Coker{\mathop{\mathrm{Coker}}}
\def\Coim{\mathop{\mathrm{Coim}}}

% Boxtimes
%
\DeclareMathSymbol{\boxtimes}{\mathbin}{AMSa}{"02}

%
% Macros for moduli stacks/spaces
%
\def\QCohstack{\mathcal{QC}\!\mathit{oh}}
\def\Cohstack{\mathcal{C}\!\mathit{oh}}
\def\Spacesstack{\mathcal{S}\!\mathit{paces}}
\def\Quotfunctor{\mathrm{Quot}}
\def\Hilbfunctor{\mathrm{Hilb}}
\def\Curvesstack{\mathcal{C}\!\mathit{urves}}
\def\Polarizedstack{\mathcal{P}\!\mathit{olarized}}
\def\Complexesstack{\mathcal{C}\!\mathit{omplexes}}
% \Pic is the operator that assigns to X its picard group, usage \Pic(X)
% \Picardstack_{X/B} denotes the Picard stack of X over B
% \Picardfunctor_{X/B} denotes the Picard functor of X over B
\def\Pic{\mathop{\mathrm{Pic}}\nolimits}
\def\Picardstack{\mathcal{P}\!\mathit{ic}}
\def\Picardfunctor{\mathrm{Pic}}
\def\Deformationcategory{\mathcal{D}\!\mathit{ef}}
\usepackage{cleveref}
\crefname{originalreference}{original reference}{original references}
\Crefname{originalreference}{Original reference}{Original references}
\makeatletter
\newlabel{spaces-topologies-definition-fpqc-covering}{{03MQ}{source}{}{sourcepointer-1}{}}
\newlabel{spaces-topologies-definition-fpqc-covering@cref}{{[originalreference][0][]03MQ}{[1][0][]source}{}{}{}}
\makeatother
\crefname{originalreference}{original reference}{original references}
\Crefname{originalreference}{Original reference}{Original references}
\makeatletter
\newlabel{spaces-properties-proposition-sheaf-fpqc}{{0APL}{source}{}{sourcepointer-2}{}}
\newlabel{spaces-properties-proposition-sheaf-fpqc@cref}{{[originalreference][0][]0APL}{[1][0][]source}{}{}{}}
\makeatother
\crefname{originalreference}{original reference}{original references}
\Crefname{originalreference}{Original reference}{Original references}
\makeatletter
\newlabel{spaces-topologies-lemma-fpqc}{{03MR}{source}{}{sourcepointer-3}{}}
\newlabel{spaces-topologies-lemma-fpqc@cref}{{[originalreference][0][]03MR}{[1][0][]source}{}{}{}}
\makeatother
\crefname{originalreference}{original reference}{original references}
\Crefname{originalreference}{Original reference}{Original references}
\makeatletter
\newlabel{spaces-topologies-lemma-recognize-fpqc-covering}{{03MS}{source}{}{sourcepointer-4}{}}
\newlabel{spaces-topologies-lemma-recognize-fpqc-covering@cref}{{[originalreference][0][]03MS}{[1][0][]source}{}{}{}}
\makeatother
\title{Representable functors are sheaves for fpqc coverings}
\author{The Stacks Project Authors}
\date{stacks-project:spaces-descent.tex:local-snapshot-20261012; GNU FDL 1.2 or later}
\begin{document}
\maketitle
\noindent\textit{Editorial note.} All delivered mathematical text is the author's own, copied verbatim from the pinned Stacks Project chapter file spaces-descent.tex; no statement or proof step was rewritten, reordered or removed.\par\medskip
\noindent\textit{Editorial note.} The proof's closing diagram chase is carried in one sentence by the author; all ingredients (uniqueness, quotient presentation) are stated in the span.\par\medskip
\noindent\textit{Allowed standard prerequisites (editorial).}\begin{itemize}
\item fpqc sheaf property for schemes. Stacks Project, Properties of Spaces, Proposition spaces-properties-proposition-sheaf-fpqc. Supplies the equalizer property for the fpqc covering \$\textbackslash{}\{U\_i \textbackslash{}to U\textbackslash{}\}\$ of schemes.
\item fpqc coverings by base change and recognition. Stacks Project, Topologies on Spaces, Lemmas spaces-topologies-lemma-fpqc and spaces-topologies-lemma-recognize-fpqc-covering. Shows \$\textbackslash{}\{U\_i \textbackslash{}to U\textbackslash{}\}\$ is an fpqc covering.
\item Algebraic spaces as sheaf quotients. Stacks Project, Descent on Spaces, discussion of \$T = U/(U \textbackslash{}times\_T U)\$ in the chapter text. Supplies the descent of \$a\$ from \$U\$ to \$T\$.
\end{itemize}\medskip
\ifdefined\hypertarget\hypertarget{sourcepointer-1}{}\fi
\noindent\textit{Source reference.} Original reference 03MQ: sources/stacks-chow/tags.txt, source lines 4758--4758. Cross-chapter standard prerequisite cited by the author inside the counted span; it is used as a pointer and its own proof is not part of this item.\par\medskip
\ifdefined\hypertarget\hypertarget{sourcepointer-2}{}\fi
\noindent\textit{Source reference.} Original reference 0APL: sources/stacks-chow/tags.txt, source lines 13479--13479. Cross-chapter standard prerequisite cited by the author inside the counted span; it is used as a pointer and its own proof is not part of this item.\par\medskip
\ifdefined\hypertarget\hypertarget{sourcepointer-3}{}\fi
\noindent\textit{Source reference.} Original reference 03MR: sources/stacks-chow/tags.txt, source lines 4759--4759. Cross-chapter standard prerequisite cited by the author inside the counted span; it is used as a pointer and its own proof is not part of this item.\par\medskip
\ifdefined\hypertarget\hypertarget{sourcepointer-4}{}\fi
\noindent\textit{Source reference.} Original reference 03MS: sources/stacks-chow/tags.txt, source lines 4760--4760. Cross-chapter standard prerequisite cited by the author inside the counted span; it is used as a pointer and its own proof is not part of this item.\par\medskip
\begin{lemma}
\label{lemma-fpqc-universal-effective-epimorphisms}
Let $S$ be a scheme. Let $\{T_i \to T\}$ be an fpqc covering of algebraic
spaces over $S$, see Topologies on Spaces, Definition
\ref{spaces-topologies-definition-fpqc-covering}.
Then given an algebraic space $B$ over $S$ the sequence
$$
\xymatrix{
\Mor_S(T, B) \ar[r] &
\prod\nolimits_i \Mor_S(T_i, B) \ar@<1ex>[r] \ar@<-1ex>[r] &
\prod\nolimits_{i, j} \Mor_S(T_i \times_T T_j, B)
}
$$
is an equalizer diagram.
In other words, every representable functor on the category of
algebraic spaces over $S$ satisfies the sheaf condition for
fpqc coverings.
\end{lemma}

\begin{proof}
We know this is true if $\{T_i \to T\}$ is an fpqc covering of
schemes, see Properties of Spaces, Proposition
\ref{spaces-properties-proposition-sheaf-fpqc}.
This is the key fact and we encourage the reader to skip the rest
of the proof which is formal. Choose a scheme $U$ and a surjective
\'etale morphism
$U \to T$. Let $U_i$ be a scheme and let $U_i \to T_i \times_T U$
be a surjective \'etale morphism. Then $\{U_i \to U\}$ is an
fpqc covering. This follows from
Topologies on Spaces, Lemmas \ref{spaces-topologies-lemma-fpqc} and
\ref{spaces-topologies-lemma-recognize-fpqc-covering}.
By the above we have the result for $\{U_i \to U\}$.

\medskip\noindent
What this means is the following: Suppose that $b_i : T_i \to B$
is a family of morphisms with
$b_i \circ \text{pr}_0 = b_j \circ \text{pr}_1$ as morphisms
$T_i \times_T T_j \to B$. Then we let $a_i : U_i \to B$ be the
composition of $U_i \to T_i$ with $b_i$. By what was said above
we find a unique morphism $a : U \to B$ such that
$a_i$ is the composition of $a$ with $U_i \to U$.
The uniqueness guarantees that $a \circ \text{pr}_0 = a \circ \text{pr}_1$
as morphisms $U \times_T U \to B$. Then since $T = U/(U \times_T U)$
as a sheaf, we find that $a$ comes from a unique morphism $b : T \to B$.
Chasing diagrams we find that $b$ is the morphism we are looking for.
\end{proof}

\end{document}
